\documentclass{../../kalkulus-ppt}

\author[Tetew]{Teosofi Hidayah Agung}
\date{\today}
\title[Kalkulus 1 - Sup B]{Matriks, Determinan, dan SPL}
\institute[Matematika ITS]{Departemen Matematika\\ Institut Teknologi Sepuluh Nopember}
\titlegraphic{{\includegraphics[scale=0.3]{ITS.png}$\quad$\includegraphics[scale=0.02]{M.png}}}

\newcommand{\dom}{\mathcal{D}}
\newcommand{\rng}{\mathcal{R}}

\renewcommand{\arraystretch}{1.5}

\begin{document}
\begin{frame}
    \titlepage
\end{frame}

\AtBeginSection{
    {
        \begin{frame}{Daftar isi}
            \tableofcontents[currentsection]
        \end{frame}}
}

\section{Matriks dan Operasi Matriks}
\begin{frame}
    \frametitle{\insertsection}
    \begin{definisi}
        \textbf{Matriks} adalah susunan bilangan berbentuk persegi panjang yang diatur dalam baris dan kolom. Bilangan-bilangan dalam matriks disebut entri atau elemen matriks.
    \end{definisi}
    \begin{block}{Operasi Matriks}
        Beberapa operasi dasar pada matriks antara lain:
        \begin{itemize}
            \item Penjumlahan matriks
            \item Perkalian skalar dengan matriks
            \item Perkalian dua matriks
            \item Transpos matriks
        \end{itemize}
    \end{block}
\end{frame}

\section{Operasi Baris Elementer dan Matriks Invers}
\begin{frame}
    \frametitle{\insertsection}
    \begin{block}{Operasi Baris Elementer (OBE)}
        Tiga operasi baris elementer pada suatu matriks:
        \begin{enumerate}
            \item Menukar posisi dua baris.
            \item Mengalikan suatu baris dengan skalar tak-nol.
            \item Menambahkan kelipatan suatu baris ke baris yang lain.
        \end{enumerate}
    \end{block}
    \onslide<2->{
        \begin{definisi}[Matriks Invers]
            Jika $A$ adalah matriks persegi, dan jika terdapat matriks $B$ yang berukuran sama sedemikian sehingga $AB = BA = I$, maka $A$ dikatakan \textbf{dapat dibalik} (invertible) dan $B$ disebut \textbf{invers} dari $A$ ($A^{-1}$).
        \end{definisi}
    }
\end{frame}

\section{Sistem Persamaan Linier}
\begin{frame}
    \frametitle{\insertsection}
    \begin{definisi}[Sistem Persamaan Linier]
        Sistem Persamaan Linier (SPL) dengan $m$ persamaan dan $n$ variabel $x_1, x_2, \dots, x_n$ dituliskan sebagai:
        \begin{align*}
            a_{11}x_1 + a_{12}x_2 + \dots + a_{1n}x_n &= b_1 \\
            a_{21}x_1 + a_{22}x_2 + \dots + a_{2n}x_n &= b_2 \\
            \vdots \qquad \qquad \qquad &= \vdots \\
            a_{m1}x_1 + a_{m2}x_2 + \dots + a_{mn}x_n &= b_m
        \end{align*}
        Dapat ditulis dalam bentuk matriks $AX = B$.
    \end{definisi}
\end{frame}

\section{Determinan}
\begin{frame}
    \frametitle{\insertsection}
    \begin{definisi}[Determinan]
        Untuk matriks persegi $A$ berukuran $2 \times 2$:
        \[A = \begin{pmatrix} a & b \\ c & d \end{pmatrix} \implies \det(A) = ad - bc\]
        Untuk matriks berukuran $n \times n$, determinan dihitung menggunakan ekspansi kofaktor atau menggunakan OBE.
    \end{definisi}
\end{frame}

\section{Nilai Eigen dan Vektor Eigen}
\begin{frame}
    \frametitle{\insertsection}
    \begin{definisi}[Nilai Eigen dan Vektor Eigen]
        Misalkan $A$ adalah matriks persegi $n \times n$. Sebuah skalar $\lambda$ disebut \textbf{nilai eigen} (eigenvalue) dari $A$ jika terdapat vektor tak-nol $\mathbf{x}$ di $\mathbb{R}^n$ sedemikian sehingga
        \[A\mathbf{x} = \lambda \mathbf{x}\]
        Vektor $\mathbf{x}$ tersebut disebut \textbf{vektor eigen} (eigenvector) dari $A$ yang berkorespondensi dengan $\lambda$.
    \end{definisi}
\end{frame}

\end{document}
